\chapter{Q Fever Antibody Test}

\section*{What Is This Test?}
Q fever antibody testing detects antibodies against \textit{Coxiella burnetii}, an intracellular bacterium transmitted mainly by inhalation of contaminated aerosols from infected livestock and birth products. Infection may cause an acute febrile illness, hepatitis, pneumonia, or chronic endocarditis and vascular infection.

\section*{Alternative Names}
Q Fever Serology; \textit{Coxiella burnetii} Antibodies; Q Fever Phase I and Phase II Antibodies; \textit{Coxiella} IFA

\section*{How It Is Performed}
Serum is tested by indirect immunofluorescence or another validated immunoassay for phase I and phase II IgG and IgM antibodies. Antibody levels are often compared between an acute specimen and a convalescent specimen collected several weeks later.

\section*{Why It Is Ordered}
\begin{itemize}
  \item Evaluation of prolonged fever, atypical pneumonia, hepatitis, or severe febrile illness after livestock or abattoir exposure
  \item Investigation of suspected acute Q fever when routine cultures are negative
  \item Assessment of chronic Q fever endocarditis or vascular infection in a high-risk patient
\end{itemize}

\section*{Interpretation and Limitations}
A single early negative result does not exclude acute Q fever because antibodies may not yet be detectable. A fourfold or greater rise in phase II antibody between paired specimens supports recent infection. High phase I IgG titres, with compatible illness and persistent infection, support chronic Q fever but are not diagnostic alone. Cross-reactivity and prior exposure can complicate interpretation; results should be reviewed with an infectious-disease specialist when chronic disease is suspected.
\section*{Regulatory Status and Authorization Date}
Regulatory status applies to the specific assay, instrument, manufacturer, and intended use rather than to the test name alone. No single approval date can be assigned to this test category. For a commercial assay, verify the current product in the relevant regulator's database: the U.S. Food and Drug Administration (FDA) clearance, approval, or authorization; India's Central Drugs Standard Control Organisation (CDSCO) licence or permission; the European Union In Vitro Diagnostic Regulation (EU IVDR) CE certificate issued through the applicable conformity-assessment route; or the United Kingdom Medicines and Healthcare products Regulatory Agency (MHRA) registration and UKCA/CE status. Laboratory-developed tests may not have an individual FDA, CDSCO, EU IVDR, or MHRA approval date.
\clearpage
